%%%%%%%%%%%%%%%%%%%%%%%%
%
% $Autor: Gautam Ramesh$
% $Datum: 2025-10-14 22:01:30Z $
% $Pfad: /Users/gautamramesh/githubNew/ML26-03-Car-Licence-Plate-Identification-with-a-Jetson-Nano/Assembly - edit/Chapters/en/Troubleshooting.tex $
% $Version: 1 $
%
% !TeX encoding = utf8
% !TeX root = Rename
%
%%%%%%%%%%%%%%%%%%%%%%%%



\chapter{Troubleshooting}

This section provides diagnostic guidance for common hardware issues that may arise during setup or operation of the Jetson Nano–based License Plate Detection System.  
It helps identify potential faults in power, display, camera, and peripheral connections to restore functionality quickly and safely.

\section{Initial Inspection}

Before powering on the system, ensure the following checks:
\begin{itemize}
	\item Verify that all cables are firmly seated in their respective ports (USB, HDMI, power barrel jack).
	\item Inspect for bent GPIO pins or loose microSD insertion.
	\item Confirm that the power adapter outputs a stable \SI{5}{\volt}/\SI{4}{\ampere}.
	\item Use a clean, static-free workspace to prevent ESD-related failures.
\end{itemize}

% -------------------------------------------------------
% Table: Common Issues and Solutions
% -------------------------------------------------------
\begin{table}[H]
	\centering
	\caption{Common Hardware Issues and Corrective Actions}
	\begin{tabular}{|L{3.8cm}|L{7.5cm}|L{2.5cm}|}
		\hline
		\textbf{Symptom} & \textbf{Possible Cause / Solution} & \textbf{Priority} \\ \hline
		No Power LED / Boot Failure & 
		\begin{itemize}[leftmargin=*]
			\item Check DC adapter output with multimeter.  
			\item Reseat the power connector and verify correct polarity.  
			\item Replace microSD card if image is corrupted.  
		\end{itemize} & High \\ \hline
		
		No HDMI Display Output & 
		\begin{itemize}[leftmargin=*]
			\item Confirm monitor input source (HDMI 1/2).  
			\item Use known-good cable and test at \texttt{1920x1080} resolution.  
			\item Reflash JetPack image if system hangs during boot logo.  
		\end{itemize} & High \\ \hline
		
		Camera Not Detected & 
		\begin{itemize}[leftmargin=*]
			\item Reconnect USB cable to another port.  
			\item Check via terminal using \texttt{lsusb} or \texttt{v4l2-ctl --list-devices}.  
			\item Replace cable or camera if not recognized.  
		\end{itemize} & Medium \\ \hline
		
		System Overheating or Throttling & 
		\begin{itemize}[leftmargin=*]
			\item Verify that heatsink and fan are unobstructed.  
			\item Maintain ambient temperature below \SI{35}{\celsius}.  
			\item Use USB-powered cooling fan for long inference sessions.  
		\end{itemize} & Medium \\ \hline
		
		USB Peripherals Unresponsive & 
		\begin{itemize}[leftmargin=*]
			\item Avoid using high-draw devices on unpowered hubs.  
			\item Test each port individually with known working devices.  
			\item Reboot and check kernel logs via \texttt{dmesg | grep usb}.  
		\end{itemize} & Medium \\ \hline
		
	\end{tabular}
\end{table}

\section{Hardware Recovery Steps}

If the system fails to boot or becomes unresponsive:
\begin{enumerate}
	\item Power off the Jetson Nano completely and disconnect all peripherals.
	\item Remove and reinsert the microSD card to ensure proper contact.
	\item Inspect the power adapter for overcurrent protection tripping.
	\item Boot with only essential components: power, HDMI, and microSD.
	\item If the board still fails, reflash the SD card using the official \texttt{JetPack~5.1.2} image.
\end{enumerate}

\section{Camera-Specific Checks}

\begin{itemize}
	\item Verify the Microsoft LifeCam Show appears under \texttt{/dev/video0}.  
	\item Run \texttt{v4l2-ctl --all} to confirm driver compatibility.  
	\item If the camera shows a black feed, check ambient lighting and lens focus.  
	\item Test using \texttt{cheese} or \texttt{fswebcam} to confirm image capture.  
\end{itemize}

\section{Preventive Maintenance Notes}

\begin{itemize}
	\item Maintain adequate ventilation around the Jetson Nano to prevent thermal stress.
	\item Avoid frequent insertion/removal of USB peripherals to reduce wear on ports.
	\item Use antistatic wrist straps when touching the board or GPIO pins.
	\item Keep firmware and models updated but only from verified NVIDIA or project repositories.
\end{itemize}

\subsection*{Summary}

Proper troubleshooting ensures minimal downtime and maximizes hardware longevity.  
Most operational faults can be resolved through systematic checks of power, connection integrity, and peripheral recognition.  
In persistent failure scenarios, the system can be safely recovered by reflashing the OS and restoring configuration backups.